\documentclass[11pt]{article}
\usepackage{amsmath,amssymb,amsthm,mathtools}
\usepackage{enumitem}
\usepackage{geometry}
\geometry{margin=1in}

\title{Step 72: Candidate Completed Weil Feature Maps for the RH Membrane Route}
\author{Working note}
\date{}

\newtheorem{theorem}{Theorem}
\newtheorem{proposition}{Proposition}
\newtheorem{lemma}{Lemma}
\newtheorem{definition}{Definition}
\newtheorem{warning}{Warning}
\newtheorem{obligation}{Obligation}
\newtheorem{remark}{Remark}

\newcommand{\K}{\mathsf K}
\newcommand{\AZ}{\mathsf A_Z}
\newcommand{\T}{\mathcal T}
\newcommand{\Wcmp}{W_{\rm cmp}}
\newcommand{\Wpr}{W_{\rm pr}}
\newcommand{\Winf}{W_{\infty}}
\newcommand{\Wpole}{W_{\rm pole}}
\newcommand{\Wtail}{W_{\rm tail}}
\newcommand{\EF}{E_{\rm EF}}
\newcommand{\supp}{\operatorname{supp}}
\newcommand{\Ran}{\operatorname{Ran}}

\begin{document}
\maketitle

\section{Purpose}
This note is the first RH-specific construction step after the abstract
formed-layer membrane framework.  The framework asks for a Douglas-type
contractive bridge
\[
  V_Z = T\Wcmp, \qquad \|T\|\le 1,
\]
where \(V_Z\) is the zero-side anti-invariant displacement readout and
\(\Wcmp\) is a positive completed carrier feature map.  Equivalently,
\[
  \AZ=V_Z^*V_Z \preceq \Wcmp^*\Wcmp=\K_{\rm cmp}.
\]
The aim here is not to prove this bridge for zeta.  The aim is to choose a
concrete test-function side and write the candidate completed Weil carrier in
feature-map language, separating the prime-power, archimedean, pole/completion,
and tail slots.

\section{Admissible test side}
Let \(\T\subset \mathcal S(\mathbb R)\) be a real, even, smooth test-function
space in the logarithmic variable \(t=\log x\).  We write
\[
  (\tau_a f)(t)=f(t+a),
  \qquad
  \widehat f(\xi)=\int_{\mathbb R} f(t)e^{-i\xi t}\,dt.
\]
The evenness condition is not essential for the algebra below, but it is the
natural starting point for the completed self-dual zeta package.

The anti-linear RH symmetry is
\[
  J(s)=1-\overline{s}.
\]
In centered coordinates \(s=1/2+x+i\gamma\), the anti-invariant coordinate is
\[
  \psi_-(s)=x=\Re(s)-\frac12.
\]
A zero-side displacement feature can therefore be represented abstractly as
\[
  (V_Z f)_\rho = \sqrt{w_\rho}\,\bigl(\Re\rho-\tfrac12\bigr)\,\mathcal E_\rho(f),
\]
where \(\mathcal E_\rho\) is the selected zero-evaluation/readout and
\(w_\rho\ge0\) is the zero-ledger weight.  This is a formal visible zero ledger;
the completed zero-side visibility and tail records are separate obligations.

\section{Positive shift-feature template}
A positive feature map has the form
\[
  (W_\mu f)(a,t)=\sqrt{\mu(a)}\,\bigl(f(t+a)-f(t)\bigr),
\]
for a positive measure or weighted counting measure \(\mu\) on shifts.  Then
\[
  \|W_\mu f\|^2
  =\int \mu(a)\,\|\tau_a f-f\|_{L^2_t}^2\,da
  =\int_{\mathbb R}\Psi_\mu(\xi)|\widehat f(\xi)|^2\,d\xi,
\]
where
\[
  \Psi_\mu(\xi)=\int \mu(a)|e^{ia\xi}-1|^2\,da\ge0.
\]
This template is useful because it turns arithmetic shifts into an honest
positive Dirichlet-form feature.  The crucial question is whether the completed
Weil form can be represented, after lawful quotienting and completion, by a sum
of such positive features plus finite-dimensional completion features.

\section{Candidate four-feature decomposition}
For a support/refinement parameter \(L\), define a smooth nonnegative cutoff
\(\alpha_L(a)\in[0,1]\), with \(\alpha_L(a)=1\) on the accepted visible range
and \(\alpha_L(a)=0\) outside the finite/refined range.  A candidate completed
feature map is
\[
  \Wcmp f
  =
  \begin{bmatrix}
    \Wpr,L f\\
    \Winf f\\
    \Wpole f\\
    \Wtail,L f
  \end{bmatrix}.
\]
The four slots are as follows.

\subsection{Prime-power slot}
A natural positive prime-power shift feature is
\[
  (\Wpr,L f)_{n}(t)
  =
  \sqrt{\frac{\Lambda(n)}{\sqrt n}}\,\alpha_L(\log n)
  \bigl(f(t+\log n)-f(t)\bigr),
\]
for \(n\ge2\).  Its quadratic form is
\[
  \|\Wpr,L f\|^2
  =
  \sum_{n\ge2}\frac{\Lambda(n)}{\sqrt n}\alpha_L(\log n)^2
  \|\tau_{\log n}f-f\|^2.
\]
This is positive by construction.  It is not, by itself, the classical prime
side of the explicit formula; rather, it is the positive feature candidate that
would have to be shown to dominate the anti-invariant zero displacement after
completion.

\subsection{Archimedean/gamma slot}
The archimedean slot is modeled by a positive continuous shift feature
\[
  (\Winf f)(a,t)
  =\sqrt{\kappa_\infty(a)}\,\bigl(f(t+a)-f(t)\bigr),
  \qquad a>0,
\]
where one formal candidate kernel is
\[
  \kappa_\infty(a)=\frac{1}{2\sinh(a/2)}.
\]
The kernel has a singularity at \(a=0\), so this slot requires a declared
renormalization or a quotient by the zero-shift/pole modes.  The positive-slot
record is therefore:
\[
  \Winf \text{ is accepted only after the archimedean singular part is
  quotient-handled or paired with }\Wpole.
\]

\subsection{Pole/completion slot}
The pole/completion feature is finite-dimensional.  It records the completion
and null/pole constraints that remove illegal zero-energy modes.  Abstractly,
\[
  \Wpole f=(\langle f,p_1\rangle,\ldots,\langle f,p_m\rangle)
\]
for declared completion functionals \(p_i\).  This slot is not optional: omitting
it can create a fake zero budget or a false tail-free explicit formula.

\subsection{Tail/support slot}
The tail slot records omitted shifts, support leakage, and promotion defects:
\[
  \Wtail,L f
  =\text{positive feature for omitted visible shifts and support/tail defects}.
\]
At a finite stage, this may be represented only by a defect matrix
\(E_{{\rm tail},L}\).  A finite-window claim is exact only if
\[
  E_{{\rm tail},L}\to0
\]
in the fixed-ledger or exhaustive-ledger sense.

\section{The candidate gate}
The Step--72 candidate gate is the following.

\begin{obligation}[Candidate Weil--Douglas feature-map gate]
Choose \(\T\), zero readout \(V_Z\), and completed features
\(\Wpr,\Winf,\Wpole,\Wtail\).  The bridge gate passes only if there exists a
contraction \(T\) such that
\[
  V_Z=T
  \begin{bmatrix}
    \Wpr\\ \Winf\\ \Wpole\\ \Wtail
  \end{bmatrix},
  \qquad \|T\|\le1.
\]
Equivalently,
\[
  \AZ\preceq
  \K_{\rm pr}+\K_\infty+\K_{\rm pole}+\K_{\rm tail}.
\]
If only a visible finite carrier has been constructed, the honest statement is
\[
  \AZ\preceq \K_{\rm pr}+\K_\infty+\K_{\rm pole}+E_{\rm tail}+E_{\rm EF}.
\]
\end{obligation}

\section{Signed explicit formula versus positive feature maps}
The classical explicit formula contains signed terms.  A signed identity of the
form
\[
  Z(f)=P_{\rm pr}(f)-P_\infty(f)+P_{\rm pole}(f)+P_{\rm tail}(f)
\]
does not automatically provide a positive feature map.  The Douglas gate needs
a Loewner domination
\[
  V_Z^*V_Z\preceq \sum_i W_i^*W_i,
\]
not a scalar equality of distributions or traces.

\begin{warning}[Trace balance is not domination]
Even if
\[
  \operatorname{tr}(V_Z^*V_Z)=\operatorname{tr}(\Wcmp^*\Wcmp),
\]
it need not follow that
\[
  V_Z^*V_Z\preceq\Wcmp^*\Wcmp.
\]
The bridge is accepted only by Loewner domination, equivalently by Douglas
factorization.
\end{warning}

\section{Classification of the four slots}
\begin{center}
\begin{tabular}{lll}
Slot & Positive-feature status & Remaining obligation\\
\hline
Prime-power & Positive as shift Dirichlet feature & Show it matches/completes prime side\\
Archimedean & Positive after renormalization/quotient & Handle singularity and gamma normalization\\
Pole/completion & Positive finite-dimensional feature & Declare pole/null quotient and completion modes\\
Tail/support & Positive or defect-budgeted & Prove vanishing tail/exhaustive ledger\\
\end{tabular}
\end{center}

\section{Nonclaim boundary}
This note does not prove RH.  It does not prove the Weil positivity criterion in
Douglas form.  It does not prove that the shift-feature candidate is the correct
zeta carrier.  It only identifies a concrete feature-map target and separates
which classical pieces are already positive candidates from which pieces require
completion, quotienting, or defect records before the Douglas gate can be
attempted.

\section{Next target}
The next target is to compare this candidate feature-map gate with an actual
Weil positivity formulation.  The question is:
\[
  \text{Does Weil positivity imply, or can it be strengthened to, the Douglas
  domination }V_Z^*V_Z\preceq\Wcmp^*\Wcmp?
\]
If not, the gap between scalar/bilinear Weil positivity and operator-valued
contractive domination must be recorded as the central RH bridge defect.

\end{document}
